\section{Conclusion}
\label{sec:conclusion}

IVAD frames agent-failure diagnosis as a selective decision over
machine-checkable causal claims rather than an unconstrained explanation. It
combines re-addressable trajectory evidence, deterministic eligibility,
controlled semantic verification, bounded revision and acquisition, and a
group-selective acceptance protocol. Under a frozen pipeline, independently
sampled preregistered groups, and simultaneous finite-candidate bounds, the
protocol provides a calibration-only PAC statement with an explicit
no-feasible-point outcome.

SpanVouch realizes this trust boundary as a versioned, recoverable open-source
system. On the frozen 36-candidate benchmark, deterministic verification
accepted all 20 valid reports and intercepted all 16 injected defects. The
audited release gate recorded 1,637 passes, one platform-specific skip, and
93.40\% coverage, while the 24-cell offline matrix reproduced the complete
multi-domain, multi-framework artifact path with zero provider and GPU calls.
These results establish executable contract behavior, system integrity, and
reproducibility. They do not turn fake-provider execution into a semantic
effectiveness result or an empirical selective-risk certificate. Together,
IVAD and SpanVouch provide a precise protocol and an engineering substrate for
auditable acceptance of agent-failure diagnoses.
